\documentclass[10pt]{article}

\usepackage[utf8]{inputenc}

\usepackage[T1]{fontenc}

\usepackage{amsmath}

\usepackage{amsfonts}

\usepackage{amssymb}

\usepackage[version=4]{mhchem}

\usepackage{stmaryrd}



\begin{document}

целое число, может иметь не более, чем конечное тисло, решений в целых числах $x$ и $y$.



Окончательную оценку величины $v$ в неравенстве (1) получіл K. Рот (K. Roth, [2]) путем обобщения Т. м. на случай многочлсна любого числа переменных, подобного многочлену $f(x, y)$, и использования больтего числа решений неравенства (1). Результат - т е о р ем а Т уә$n \geqslant 2$. Т. м. имеет обобщение на случай приближения алгебраич. чиселI алгебраич. числами. 'Т. м. является общим методом доказательства конечности целых точек на широкого класса кривых алгебраич. многообразий (см. Диобантова геометрия, Диобантово множество). Наряду с әтим Т. м. имеет существенный недостаток является неәфффективным методом в том смысле, что не дает ответа на вопрос, существуют ли на самом деле пспользуемые в доказательствах решения неравенств (1) или соответствующих уравнений (2). Таким образом, Т. м., решая вопрос о конечности числа решений уравнения (2), не дает возможности определить, разрешимо ли конкретное уравнение такого типа и каковы количествен, оценки решений $x, y$ в зависимости от $F$.



См. также Диобантовых приближений проблемы эфдективизаиии.



Jum.: [1] T h u e A., "J. reine angew. Math.», 1909, Bd 135, S. 284-305; [2] Р от К. Ф.. "Математика", 1957, Т 1, в. 1,

c. 3-18; [3] Проблемы теории диофантовых приближений, пер. с англ., М., 1974. А. Ф. Лаврик.



ТУЭ ПОЛУСИСТЕМА, п 0 л У - Т у ә с и с т е м а, с и с те м а под с т а н о в о к-см. Туэ система. ТУЭ СИСТЕМА - ассочиативное исчисление, названное по имени А. Туә, к-рый впервые сформулировал проблему распознавания равенства слов в ассоциативных системах (п р о б л е м а Т у э, см. [1]). Если при задании Т. с. допустимыми подстановками считать только подстановки правых частей соотношений вместо левых частей (т. е. исключить обратные подстановки), то получим пол у с ис тем ы Т у в (пол уТу с системы, или сис темы под т танов о к), к-рые фактически совпадают также с локальными канонич. системами Поста. Каждая Т. с. может рассматриваться как полусистема Туә, но обратное неверно.



Jum.: [1] T h u e A., "Videnskapsselskapets Skrifter. Mat. Naturv. K1.», 1914, № 10; [2] М а л ь ц е в А. И., Алгоритмы и рекурсивные функци, М., 1965 . TEOPEMА: если Адалгебраич. иррациональность, $\delta>0$ и сколь угодно мало, то существует лишь конечное число целых решений $p$ и $q>0$ ( $p$ и $q$ взаимно просты) неравенства



\[

\left|\alpha-\frac{p}{q}\right|<q^{-2-\delta} \text {. }

\]



Эта теорема является наилучшей в своем роде - число 2 в показателе степени уменьшить нельзя. Т. $-3 .-\mathrm{P}_{*}$ т. есть усиление теоремы Лиувилля (см. Лиувилля число). Результат Лиувилля последовательно усиливали А. Туэ [1], К. Зигель [2] и, наконец, К. Рот [3]. А. Туэ доказал, что если $\alpha$ - алгебраич. число степени $n \geqslant 3$, то неравенство



\[

\left|\alpha-\frac{p}{q}\right|<q^{-v}

\]



имеет лишь конечное число целых решений $p$ и $q>0$ ( $p$ и $q$ взаимно просты) при $v>\frac{n}{2}+1$. К. Зигель установил, что утверждение теоремы Туэ справедливо при $v>2 \sqrt{n}$. Окончательный вариант теоремы, сформулированный выше, получен $\mathfrak{K}$. Ротом. Имеется $p$ адический аналог Т.-3.-Р. т. Перечисленные результаты доказываются неәффективными методами (см. Диофантовых приближений проблемы әффективизации). S. I-34; [2] S i e g e l C.., "Math. Z.», 1921, Bd 10,'s. 173-213; [3] R oth K. F., "Mathematika", 1955, v. 2, № 1, p. 1-20; [4] M a h 1 e r K., Lectures on Diophantine approximations, pt 1 , Is. 1.], $1961,[5] \mathrm{R}$ i d o u t D , "Mathematika», 1958, v. 5, № 9, p. $40-48,[6] \Gamma$ е ль ф он д А. О., Трансцендентные б алгейраичсские числа, М., $1952 . \quad$ C. В. Котов.



ТЬЮРИНГА МАІИНА - название, закрепившееся за вычислительными машинали абстрактными некрого точно охарактеризованного типа. Концепция такого рода машины возникла в середине $30-\mathbf{x}$ гг. 20 в. у А. М. Тьюринга [1] в результате произведенного им анализа действий человека, выполняющего в соответствии с заранее разработанным планом те или иные вычисления, т. е. последовательные преобразования знаковых комплексов. Анализ этот, в свою очередь, был осуцествлен им с целью ретения назревшей к тому времени проблемы поиска точного математич. әквивалента для общего интуитивного представления об $а$ плгорипме. В ходе развития алгоритмов теории появился ряд модификаций первоначального тьюринговского определения. Здесь дается версия, восходящая к Э. Посту [2], - в таком виде определение Т. м. получило весьма большое раслространение (детально Т. м. описаны, напр., в [3] и [41).



Т. м. удобно представлять себе в виде автоматически функционирующего устройства, способного находиться в конечном числе в н у т р е н и и с о с то я н й й и снабженного бесконечной внешней памятью - Ј е нт о й. Среди состояний имеется два выделенных на ч а льн о е и 3 а к л ю и т е льно е. Лента разделена на клетки и неограниченна влево и вправо. В каждой клетке ленты может быть записана любая из букв нек-рого алфавита $A$ (ради единообразия удобно считать, что в пустой клетке записана «пустая буква»). В каждый момент дискретного времени Т. м. находится в одном из своих состояний и, рассматривая одну из клеток своей ленты, воспринимает записанный в ней символ - букву алфавита $A$ или пустую букву.



Находясь в какой-либо момент времени в незаключительном состоянии, Т. м. совершает ш а г, к-рый вполне определяется ее текущим состоянием и символом, воспринимаемым ею в данный момент на ленте. IIаг әтот заключается в том, что: 1) в рассматриваемой клетке записывается новый символ, быть может, совшадающий со старым или пустой; 2) мапина переходит в новое состояние, быть может, совпадающее со старым или заключительное; 3) лента сдвигается влево или вправо на одну клетку или остается на месте. Перечисление всех возможных шагов Т. м. в зависимости от текущей комбинации «незаключительное состояние + воспринимаемый символ», представленное, напр., в виде специальной таблицы с двумя входами, наз. п р ог р а м м о й, или с х е м о й, данной Т. м. В клетках этой таблицы стоят коды соответствующих шагов машины - ее к о м а н д ы. Программа Т. м. является объектом с точной структурой, и можно условиться отождествлять Т. м. с ее программой. Стремясь подчеркнуть связь так огисанной Т. м. с алфавитом $A$, обычно говорят, что әта машина есть 'Г. м. в а л ф ав и т е $A$.



Текущее полное описание Т. м. дается ее к о н ф иг у р а ц и е й, к-рая состоит из указания для данного момента следующей информации: 1) конкретного заполнения клеток ленты символами; 2) клетки, находящейся в поле зрения машины; 3) внутреннего состояния, в к-ром машина находится. Конфигурация, соответствующая заключительному состоянию Т. м., также наз. з а к л ю ч и т е л ь н о Й.



Если зафиксировать какую-либо незаключительную конфигурацию Т. м. в качестве исходной, то работа әтой машины будет заключаться в последовательном (шаг за пагом) преобразовании исходной конфигурации в соответствии с программой машины до тех пор, пока не будет достигнута заключительная конфиру-





\end{document}